% Pursuing Stacks (cotes 134-2 to 134-8): archive cote 134-7, site page 51 (PDF page 52), typescript page 493.
% typescript. Draft compared: ps6.tex:2507-2578.
% AI-transcribed by Claude Opus 5.5 (claude-opus-5-5, adaptive thinking, effort high), Message Batches API, batch ps4 (msgbatch_012Aoc1NuskATPsyBeKSYK7J), 2026-11-12; not yet checked by eye.
\scanpage{134-7/51}
\mspage{493}

does'nt.  When we wrote down a formula such as (21) on page 487,  we were
implicitly making use of the assumption that for k-modules M,N we have
a canonical isomorphism
\[ (1) \qquad \mathrm{RHom}_{k}(M,N) \hand{\xleftarrow{\sim}} \mathrm{RHom}_{\underline{O}_k}(W(M),W(N)) \ , \]
at any rate such a formula is needed if we want to view (21) as a more
explicit way of writing (20) or (20').  Using a free resolution L. of M,
and the fact that
\[ \mathrm{Ext}^{i}_{\underline{O}_k}(\underline{O}_k, \deleted{M} W(N)) \hand{\simeq} H^{i}(\mathrm{Spec}(k), W(N)) \ = 0 \ \text{if} \ \ i \hand{>} 0, \]
stemming from the cohomology properties of flat descent,  we easily get
a map (1),  but I did not check that this map is an isomorphism,  the dif-%
ficulty coming  from the fact that  W(L.)  need not be a resolution of
W(M),  unless M is flat.  This perplexity already arises in the fppf
context  -  surely Larry Breen should know the answer.  For what we are
\deleted{\unsure{reall|realt}} after here \deleted{though} it does'nt seem to  matter too much,  as the com-%
putations of loc. cit.  were of interest mainly (maybe exclusively) in the
case when M is projective or at any rate flat,  i.e. when working with
\uline{flat} unipotent bundles   -  in which case (1) is indeed an isomorphism.

The interpretation of polynomial maps between k-modules M,N  in terms
of the topological augmented coalgebras  $\Gamma_k{}^{\wedge}$ associated to these  may seem
a little forbidding to some readers.  In the all-important and typical
case when M \xout{is} \inserted[typed]{and N are} projective and of finite type,  things become quite evident
though,  by just dualizing the more familiar concepts around polynomial
functions and homomorphisms between rings of such.  Thus W(M) is just a
usual vector-bundle,\xout{\unsure{t}}hence also a\xout{n} true honest affine scheme over k,
whose affine ring is the ring of polynomial functions on M,  which can
be identified with  $\over{\mathrm{Sym}}{\mathrm{Sym}}^{\hand{*}}_{\hand{k}}(M^{\vee})$,  the symmetric algebra on the dual module :
\[ (2) \qquad W(M) \hand{\simeq} \mathrm{Spec}(\mathrm{Sym}^{\hand{*}}_{\hand{k}}(M\ )) \ , \]
and similarly for W(N).  Polynomial maps from M to N i.e. maps from the
k-scheme W(M) to the k-scheme W(N),   just correspond to k-algebra homomor-%
phisms
\[ \hand{(3)} \qquad \mathrm{Sym}^{*}_{\hand{k}}(N) \hand{\longrightarrow} \mathrm{Sym}^{*}_{\hand{k}}(M) \]
(irrespective of the graded structures).  As each of these algebras\inserted{, as a k-module,} is
a filtering direct limit of \inserted{it's} projective \inserted{sub}modules of finite type such as
\[ \hand{(4)} \qquad \mathrm{Sym}_{k}(M)(i) \ = \ \hand{\bigoplus}_{j \hand{\leq} i} \mathrm{Sym}_{k}{}^{j}(M) \qquad , \]

\margin{51}
